\sectioncentered*{Введение}
\addcontentsline{toc}{section}{Введение}

Сфера розничной торговли автомобильными запасными частями и компонентами отличается обширной номенклатурой и постоянным спросом. В условиях высокой конкуренции ручная обработка заявок и разрозненный складской учет приводят к задержкам подбора деталей, пересортице и ошибкам резервирования. Возникает объективная необходимость во внедрении автоматизированной информационной системы, обеспечивающей централизованный учет номенклатуры, контроль остатков и сквозное сопровождение клиентских заказов.

Управление каталогом автозапчастей и корзиной покупателя создает непрерывный поток запросов с поиском и фильтрацией по маркам, моделям и категориям. Надежное функционирование системы требует применения реляционной СУБД. Реляционная модель гарантирует ссылочную целостность между учетными записями, товарами, заказами и платежами. Оформление покупки опирается на строгую транзакционную изоляцию: списание складских остатков, фиксация отпускных цен и сохранение данных доставки должны выполняться атомарно, исключая коллизии при одновременном обращении клиентов.

Цель курсового проекта~-- проектирование и разработка автоматизированной информационной системы управления интернет-магазином автозапчастей, обеспечивающей учет номенклатуры, клиентскую корзину, транзакционную обработку заказов, многоформатный обмен данными и аналитическую отчетность.

Для достижения цели поставлены следующие задачи:
\begin{itemize}
	\item исследовать предметную область торговли автозапчастями и провести анализ систем-аналогов;
	\item спроектировать концептуальную, логическую и физическую схемы базы данных;
	\item разработать серверные модули работы с СУБД, обработки заказов, обмена данными и генерации отчетов;
	\item реализовать пользовательский веб-интерфейс и составить руководство пользователя и администратора.
\end{itemize}
